\chapter{Implementaci\'on y Desarrollo}\label{cap_implementacion}
% tfg-floatbarrier-auto
\makeatletter
\@ifundefined{tfgfloatbarrier}{
  \def\tfgfloatbarrier{1}
  \let\tfgorigfigure\figure
  \let\tfgendorigfigure\endfigure
  \renewenvironment{figure}{\tfgorigfigure}{\tfgendorigfigure\FloatBarrier}
  \let\tfgorigtable\table
  \let\tfgendorigtable\endtable
  \renewenvironment{table}{\tfgorigtable}{\tfgendorigtable\FloatBarrier}
}{}
\makeatother

\section{Implementaci\'on backend}

La implementaci\'on backend se estructura en el paquete principal de la aplicaci\'on, donde destacan \textbf{controladores} como entrada HTTP y validaci\'on de contratos, \textbf{servicios} para l\'ogica de negocio y reglas del dominio, \textbf{repositorios} con acceso a datos mediante Spring Data JPA y un bloque de \textbf{seguridad} con configuraci\'on central y filtros especializados. Esta separaci\'on refuerza la testabilidad y el control de responsabilidades.

Un punto clave es \texttt{EntregaService}, que centraliza validaciones de entrega,
control de versiones, restricciones temporales y operativa de evaluaci\'on,
con apoyo de servicios auxiliares para materiales y notificaciones.

\section{Implementaci\'on de flujos cr\'iticos}

Los flujos con mayor impacto funcional son la \textbf{autenticaci\'on y sesi\'on} (login, emisi\'on de tokens, refresco y cierre seguro), la \textbf{publicaci\'on de entregables} (creaci\'on por docente y configuraci\'on de fechas y condiciones), la \textbf{entrega de material} (subida de ficheros, verificaciones, versionado y almacenamiento de metadatos) y la \textbf{evaluaci\'on} (correcci\'on por profesor autorizado, feedback y visibilidad de nota). Esta selecci\'on concentra los riesgos funcionales del sistema.

Cada flujo incorpora validaciones de entrada, control de errores y verificaci\'on de permisos,
asegurando consistencia en escenarios concurrentes y previniendo acciones fuera de alcance.

\section{Implementaci\'on frontend}

La implementaci\'on del cliente cubre pantallas y operaciones necesarias para el MVP, incluyendo acceso y gesti\'on de sesi\'on, paneles de actividades y entregables, formularios de entrega y seguimiento, y vista de evaluaciones con retroalimentaci\'on. Esta cobertura garantiza recorrido completo por rol.

Se han introducido medidas pr\'acticas para mejorar usabilidad:
mensajes de validaci\'on temprana, estados visuales de proceso,
y comportamiento resiliente ante errores de red.

\section{Persistencia, migraciones e integraciones}

Se utiliza PostgreSQL en escenarios de despliegue y un perfil de desarrollo
local para iteraci\'on r\'apida. Flyway gestiona la evoluci\'on del esquema, incluyendo ampliaciones como atributos para validaci\'on estructural de entregas ZIP, campos para integraci\'on cloud (OneDrive/Cloudinary), tablas y preferencias de notificaci\'on, y tokens de recuperaci\'on de contrase\~na. Esta evoluci\'on versionada protege la consistencia hist\'orica de datos.

Las integraciones externas se dise\~nan por configuraci\'on,
de modo que el sistema puede operar en local sin dependencia obligatoria
de servicios de terceros.
En producci\'on el env\'io de correos se realiza mediante Brevo a trav\'es
de peticiones HTTPS, evitando depender del protocolo SMTP, que puede estar
restringido en plataformas PaaS. Este canal se utiliza para recuperaci\'on
de contrase\~na y notificaciones transaccionales.

\section{Verificaci\'on de implementaci\'on}

La verificaci\'on se realiza en capas (unitaria, integraci\'on, E2E y mutaci\'on),
apoyada por pipelines CI/CD y controles de seguridad en preproducci\'on.
Este enfoque proporciona evidencia objetiva de robustez y reduce el riesgo de regresiones.

\FloatBarrier
\section{Discusi\'on t\'ecnica}

\section{Decisiones y compromisos}

Las decisiones principales del dise\~no implican compensaciones expl\'icitas. La separaci\'on SPA + API mejora mantenibilidad y escalado, a costa de mayor esfuerzo de orquestaci\'on y pruebas de integraci\'on. El uso de JWT con refresh token simplifica backend stateless, pero exige controles cuidadosos de expiraci\'on y renovaci\'on. La adopci\'on de DTOs y capas reduce acoplamiento con mayor c\'odigo estructural, y las migraciones versionadas aportan trazabilidad de datos requiriendo disciplina de despliegue entre entornos.

\section{Limitaciones y trabajo futuro}

Aunque el n\'ucleo funcional est\'a cubierto, se identifican l\'ineas de evoluci\'on como la ampliaci\'on de autenticaci\'on avanzada (p. ej., 2FA completa), una observabilidad m\'as profunda con m\'etricas operativas y el cierre progresivo de requisitos \emph{Should/Could} pendientes. Estas mejoras no condicionan la validez del TFG, pero gu\'ian la continuidad t\'ecnica.

Estas mejoras no condicionan la validez del TFG,
sino que delimitan una hoja de ruta t\'ecnica razonable para continuidad del proyecto.

\FloatBarrier
\section{Conclusiones del cap\'itulo}

El an\'alisis realizado demuestra coherencia entre lo especificado y lo implementado.
Los requisitos funcionales priorizados se traducen en flujos operativos reales,
los no funcionales se abordan mediante decisiones t\'ecnicas concretas,
y el dise\~no modular facilita mantenimiento, verificaci\'on y evoluci\'on.

Por tanto, el resultado no es solo una implementaci\'on funcional,
sino una soluci\'on con fundamento de ingenier\'ia de software,
alineada con los objetivos acad\'emicos y profesionales del TFG.

\section{Justificaci\'on detallada de decisiones}\label{justificacion}

\FloatBarrier
\subsection{Introducci\'on}

En un TFG de ingenier\'ia de software, implementar funcionalidades no es suficiente
si no se explica por qu\'e se ha elegido cada opci\'on t\'ecnica,
qu\'e alternativas se consideraron y qu\'e compromisos se asumieron.

Este cap\'itulo ampl\'ia de forma exhaustiva la justificaci\'on de decisiones
adoptadas durante el proyecto,
con foco en cuatro ejes: producto, arquitectura, calidad y operaci\'on.

\FloatBarrier
\subsection{Marco de toma de decisiones}

\subsubsection{Criterios utilizados}

La selecci\'on de alternativas se realiz\'o usando un enfoque multicriterio, priorizando la alineaci\'on con requisitos funcionales y no funcionales del ERS, la viabilidad temporal dentro de 660 horas de proyecto, la mantenibilidad y capacidad de evoluci\'on posterior, el coste operativo compatible con entorno acad\'emico y la seguridad con trazabilidad del ciclo completo. Esta pauta evita decisiones aisladas y garantiza coherencia transversal.

\subsubsection{Ponderaci\'on de criterios}

Para evitar decisiones ad-hoc,
se aplic\'o una ponderaci\'on orientativa de criterios
que se mantuvo estable en los principales bloques de dise\~no.

\begin{table}[!htbp]
\centering
\begin{tabular}{|P{6.4cm}|P{2.5cm}|}
\hline
\textbf{Criterio} & \textbf{Peso (\%)} \\
\hline
Adecuaci\'on funcional y cobertura de requisitos & 30 \\
\hline
Mantenibilidad y evoluci\'on & 20 \\
\hline
Tiempo de implementaci\'on y riesgo de entrega & 20 \\
\hline
Coste operativo y sostenibilidad & 15 \\
\hline
Seguridad y robustez t\'ecnica & 15 \\
\hline
\textbf{Total} & \textbf{100} \\
\hline
\end{tabular}
\caption{Ponderaci\'on orientativa para selecci\'on de decisiones}
\label{tab:pesos-criterios}
\end{table}

\FloatBarrier
\subsection{Decisiones de alcance funcional}

\subsubsection{Decisi\'on D1: definir un MVP centrado en flujo acad\'emico nuclear}

\textbf{Alternativas consideradas}: incorporar desde inicio todas las funcionalidades del backlog o limitar el alcance a un MVP con ciclo completo de entrega y evaluaci\'on. La comparativa se centr\'o en riesgo de entrega y capacidad de validaci\'on.

\textbf{Decisi\'on adoptada}: segunda alternativa.

\textbf{Raz\'on principal}:
maximizar valor entregado en plazo,
evitando dispersi\'on en funcionalidades accesorias
que pod\'ian comprometer la estabilidad del n\'ucleo.

\textbf{Impacto positivo}: mayor probabilidad de cierre funcional completo, mejor cobertura de pruebas sobre procesos cr\'iticos y documentaci\'on m\'as s\'olida de extremo a extremo. Este impacto refuerza la validez acad\'emica del resultado.

\textbf{Compromiso asumido}:
funcionalidades \emph{Should/Could} pasan a hoja de ruta,
sin bloquear la validez del resultado principal.

\subsubsection{Decisi\'on D2: dise\~nar experiencia por rol desde el inicio}

\textbf{Decisi\'on adoptada}:
separar claramente experiencias de administrador, profesor y estudiante.

\textbf{Razones}: reduce ambig\"uedad de permisos en interfaz y API, simplifica pruebas de autorizaci\'on y casos negativos, y mejora usabilidad al mostrar solo operaciones pertinentes. Esta separaci\'on evita errores de contexto en operaciones sensibles.

\FloatBarrier
\subsection{Decisiones de arquitectura}

\subsubsection{Decisi\'on D3: arquitectura SPA + API desacoplada}

\textbf{Alternativas}: aplicaci\'on monol\'itica renderizada en servidor o frontend SPA con backend REST desacoplado. La decisi\'on se evalu\'o en t\'erminos de evoluci\'on y coste de integraci\'on.

\textbf{Decisi\'on adoptada}: SPA + API.

\textbf{Razones de peso}: separaci\'on de responsabilidades y menor acoplamiento, evoluci\'on independiente de capa de presentaci\'on y l\'ogica, y mejor ajuste a un dominio con flujos ricos de interfaz (evaluaciones, filtros, historial). Esta elecci\'on maximiza mantenibilidad en el medio plazo.

\textbf{Trade-off}:
aumenta esfuerzo de integraci\'on y de pruebas E2E,
compensado por mayor mantenibilidad a medio plazo.

\subsubsection{Decisi\'on D4: arquitectura backend por capas}

\textbf{Decisi\'on adoptada}:
patr\'on \texttt{controller-service-repository-entity} con DTOs de frontera.

\textbf{Razones}: control claro de responsabilidades, testabilidad por m\'odulo, evitaci\'on de exponer modelo de persistencia en contratos externos y facilidad para versionado de API y evoluci\'on de dominio. El patr\'on elegido reduce riesgos de acoplamiento a futuro.

\subsubsection{Decisi\'on D5: uso de DTOs y mapper dedicado}

\textbf{Problema que resuelve}:
evitar dependencia directa entre entidades JPA y respuestas HTTP.

\textbf{Razones}: minimiza acoplamiento entre capa de datos y capa API, reduce riesgo de sobreexposici\'on de atributos sensibles y permite modelar respuestas orientadas al caso de uso. Esto asegura contratos m\'as estables y seguros.

\textbf{Coste asumido}:
mayor volumen de c\'odigo estructural,
aceptado por mejora en mantenibilidad y seguridad de contrato.

\FloatBarrier
\subsection{Decisiones de stack tecnol\'ogico}

\subsubsection{Decisi\'on D6: backend en Spring Boot + Java 21}

\textbf{Razones principales}: ecosistema maduro en seguridad, validaci\'on y persistencia, integraci\'on robusta con Spring Security, JPA y testing, y soporte industrial con documentaci\'on extensa. Esta base reduce incertidumbre en implementaci\'on.

\textbf{Comparativa resumida}:
frente a opciones como Node.js o Django,
Spring reduce incertidumbre de ensamblaje en aspectos cr\'iticos
(seguridad por rol, transacciones, control de acceso por contexto).

\subsubsection{Decisi\'on D7: frontend en React + TypeScript + Vite}

\textbf{Razones}: equilibrio entre productividad y control tipado, ecosistema amplio para rutas, formularios y testing, y arranque y build r\'apidos con Vite. Esta combinaci\'on facilita entregas incrementales sin perder control de calidad.

\textbf{Compromiso}:
mayor libertad implica mayor disciplina de arquitectura en frontend,
mitigada mediante organizaci\'on por m\'odulos (pages/components/context/services).

\FloatBarrier
\subsection{Decisiones de seguridad}

\subsubsection{Decisi\'on D8: autenticaci\'on JWT con refresh token}

\textbf{Motivaci\'on}:
compatible con arquitectura desacoplada y despliegue stateless.

\textbf{Razones de elecci\'on}: reduce estado de sesi\'on en servidor, habilita renovaci\'on transparente de sesi\'on en cliente y facilita escalado horizontal del backend. Esto mejora la coherencia con un modelo API-first.

\textbf{Riesgos identificados}:
control de expiraci\'on y revocaci\'on m\'as complejo que sesi\'on tradicional.

\textbf{Mitigaci\'on aplicada}:
expiraciones diferenciadas, refresh controlado,
filtros de seguridad y validaci\'on de contexto por recurso.

\subsubsection{Decisi\'on D9: autorizaci\'on por rol + contexto acad\'emico}

\textbf{Raz\'on}:
un rol global no es suficiente en un dominio donde un usuario
puede operar con permisos distintos seg\'un curso, grupo o actividad.

\textbf{Aplicaci\'on concreta}:
validaci\'on en controladores y servicios de pertenencia real
al contexto objetivo antes de permitir operaciones sensibles.

\FloatBarrier
\subsection{Decisiones de datos e integraciones}

\subsubsection{Decisi\'on D10: PostgreSQL como base de datos principal}

\textbf{Razones}: robustez transaccional y consistencia, buena compatibilidad con JPA/Hibernate y adaptaci\'on sencilla a entornos gestionados. Esta elecci\'on reduce riesgos de despliegue.

\subsubsection{Decisi\'on D11: Flyway como autoridad de esquema}

\textbf{Razones}: trazabilidad expl\'icita de cambios en base de datos, despliegues reproducibles entre entornos y menor riesgo de divergencias silenciosas de esquema. Se prioriza control sobre automatismos impl\'icitos.

\textbf{Decisi\'on asociada}:
usar \texttt{ddl-auto=validate} en producci\'on para evitar
mutaciones impl\'icitas del esquema fuera de control.

\subsubsection{Decisi\'on D12: estrategia h\'ibrida de almacenamiento de archivos}

\textbf{Decisi\'on adoptada}:
almacenamiento local en desarrollo y cloud persistente en producci\'on.

\textbf{Razones}: simplicidad y velocidad en entorno local, persistencia y resiliencia en entornos con filesystem ef\'imero, y capacidad de habilitar integraciones externas sin romper el n\'ucleo. Se busca equilibrio entre agilidad y robustez operativa.

\FloatBarrier
\subsection{Decisiones de calidad y verificaci\'on}

\subsubsection{Decisi\'on D13: testing en capas}

\textbf{Razones}: unitarias para retroalimentaci\'on r\'apida, integraci\'on para contratos y seguridad, E2E para validar experiencia real por rol y mutaci\'on para medir eficacia real de la suite. Esta mezcla permite cubrir profundidad y visi\'on de conjunto.

\textbf{Resultado esperado}:
combinar profundidad y coste de mantenimiento,
evitar dependencia exclusiva de una sola capa de pruebas.

\subsubsection{Decisi\'on D14: incluir mutaci\'on (PIT/Stryker)}

\textbf{Raz\'on de fondo}:
la cobertura porcentual no garantiza detecci\'on efectiva de fallos.

\textbf{Valor aportado}:
los mutantes supervivientes ayudan a localizar huecos de aserciones
y mejorar calidad de tests en c\'odigo cr\'itico.

\FloatBarrier
\subsection{Decisiones de operaci\'on y despliegue}

\subsubsection{Decisi\'on D15: GitHub Actions como columna vertebral CI/CD}

\textbf{Razones}: integraci\'on nativa con repositorio, trazabilidad de checks por commit/PR y menor sobrecarga de operaci\'on para equipo reducido. Esta decisi\'on minimiza fricci\'on en el ciclo diario.

\subsubsection{Decisi\'on D16: despliegue en PaaS gestionado}

\textbf{Razones}: minimiza administraci\'on de infraestructura, acelera iteraci\'on y validaci\'on de cambios y mantiene coste compatible con contexto acad\'emico. Se prioriza estabilidad operativa frente a control extremo.

En la implementaci\'on concreta se utiliza Render como PaaS gestionado para
el despliegue de los servicios y PostgreSQL gestionado. Para el canal de
correo se adopta Brevo mediante API HTTP, evitando depender de SMTP en
producci\'on, ya que algunos PaaS restringen puertos salientes asociados al
protocolo de correo. Inicialmente se utiliz\'o SendGrid como proveedor
de correo transaccional; tras la finalizaci\'on de su plan gratuito
se migr\'o a Brevo, que permite enviar correo a cualquier destinatario
en su plan gratuito.
El dise\~no adaptativo del \texttt{EmailService} soporta m\'ultiples
proveedores (Brevo, Resend, SendGrid) mediante un patr\'on desacoplado
que permite migrar de proveedor sin modificar l\'ogica de negocio.

\textbf{Trade-off}:
menor control fino de infraestructura frente a VPS/Kubernetes,
considerado aceptable para objetivos y alcance del TFG.

\FloatBarrier
\subsection{Matriz consolidada de decisiones}

\begin{table}[!htbp]
\centering
\begin{tabular}{|P{1.6cm}|P{3.8cm}|P{3.2cm}|P{4.4cm}|}
\hline
\textbf{ID} & \textbf{Decisi\'on} & \textbf{Alternativa principal descartada} & \textbf{Raz\'on dominante} \\
\hline
D1 & Alcance MVP priorizado & Backlog completo desde inicio & Reducir riesgo de no entrega funcional \\
\hline
D3 & SPA + API desacoplada & Monolito SSR & Mantenibilidad y evoluci\'on de UX \\
\hline
D6 & Spring Boot + Java & Node.js/Django & Menor incertidumbre en seguridad/persistencia \\
\hline
D7 & React + TS + Vite & Angular/Vue & Equilibrio productividad-control tipado \\
\hline
D8 & JWT + refresh & Sesi\'on server-side & Ajuste natural a cliente API-first \\
\hline
D11 & Flyway + validate & ddl-auto update & Trazabilidad y control de esquema \\
\hline
D12 & Cloud en producci\'on & Solo almacenamiento local & Evitar p\'erdida en entornos ef\'imeros \\
\hline
D13 & Testing en capas + mutaci\'on & Solo unitarias/cobertura & Robustez real frente a regresiones \\
\hline
D16 & PaaS gestionado & VPS/Kubernetes & Tiempo y coste operativos contenidos \\
\hline
\end{tabular}
\caption{Resumen de decisiones clave y justificaci\'on dominante}
\label{tab:resumen-decisiones}
\end{table}

\FloatBarrier
\subsection{An\'alisis coste-beneficio por bloques de decisiones}

Para reforzar la justificaci\'on,
se analiza cada bloque de decisiones en t\'erminos de coste inicial,
beneficio obtenido y retorno en el contexto del TFG.

\subsubsection{Bloque de arquitectura}

\textbf{Coste inicial}:
incremento de complejidad de integraci\'on entre frontend y backend,
definici\'on de contratos API y mayor superficie de pruebas.

\textbf{Beneficio}:
separaci\'on de responsabilidades,
facilidad para aislar incidencias,
evoluci\'on independiente de interfaces sin rehacer l\'ogica de negocio.

\textbf{Retorno}:
alto. El coste adicional de arranque se recupera
al reducir deuda t\'ecnica estructural y facilitar extensiones del sistema.

\subsubsection{Bloque de stack tecnol\'ogico}

\textbf{Coste inicial}:
tiempo de configuraci\'on de toolchain,
definici\'on de convenciones de proyecto y aprendizaje cruzado entre capas.

\textbf{Beneficio}:
ecosistema maduro,
documentaci\'on abundante,
integraci\'on directa con seguridad, testing y despliegue automatizado.

\textbf{Retorno}:
alto. La reducci\'on de incertidumbre durante la implementaci\'on
ha compensado sobradamente la configuraci\'on inicial.

\subsubsection{Bloque de seguridad}

\textbf{Coste inicial}:
dise\~no de flujo de tokens,
gesti\'on de expiraciones y tratamiento de errores de autenticaci\'on.

\textbf{Beneficio}:
control de acceso robusto,
alineaci\'on con arquitectura stateless,
capacidad de endurecimiento incremental.

\textbf{Retorno}:
muy alto. En un sistema acad\'emico con datos sensibles,
la inversi\'on en seguridad es estructural, no opcional.

\subsubsection{Bloque de calidad y pruebas}

\textbf{Coste inicial}:
esfuerzo significativo en dise\~no de suites,
mantenimiento de mocks y pipelines de verificaci\'on.

\textbf{Beneficio}:
disminuci\'on de regresiones,
detenci\'on temprana de defectos,
confianza objetiva para despliegue continuo.

\textbf{Retorno}:
muy alto. La calidad de salida mejora de forma medible
y se reduce el tiempo de correcci\'on en fases tard\'ias.

\subsubsection{Bloque de despliegue y operaci\'on}

\textbf{Coste inicial}:
configuraci\'on de entornos,
variables seguras y automatismos de publicaci\'on.

\textbf{Beneficio}:
repetibilidad,
trazabilidad de releases,
validaci\'on continua sobre entorno similar a producci\'on.

\textbf{Retorno}:
alto. La operativa se vuelve estable y menos dependiente
de ejecuciones manuales fr\'agiles.

\FloatBarrier
\subsection{Riesgos asociados a decisiones y medidas de contingencia}

Ninguna decisi\'on de ingenier\'ia es neutra;
por ello se explicitan riesgos y mitigaciones por cada eje principal.

\begin{table}[!htbp]
\centering
\begin{tabular}{|P{3.4cm}|P{4.3cm}|P{4.5cm}|}
\hline
\textbf{Eje de decisi\'on} & \textbf{Riesgo principal} & \textbf{Mitigaci\'on aplicada} \\
\hline
Arquitectura desacoplada & Errores de contrato entre cliente y API & DTOs expl\'icitos, pruebas de integraci\'on y E2E \\
\hline
JWT + refresh & Gesti\'on incorrecta de expiraciones/revocaci\'on & Pol\'iticas de expiraci\'on, filtros dedicados y validaci\'on de contexto \\
\hline
Flyway como autoridad & Fallos de migraci\'on en despliegue & Versionado incremental, validaci\'on previa y backups de datos \\
\hline
Cloud en producci\'on & Dependencia de servicio externo & Configuraci\'on por entorno y estrategia fallback local \\
\hline
CI/CD automatizada & Falsos positivos por inestabilidad de tests & Separaci\'on por etapas, reintentos controlados y revisi\'on de flaky tests \\
\hline
\end{tabular}
\caption{Riesgos derivados de decisiones y medidas de contingencia}
\label{tab:riesgos-decisiones}
\end{table}

\FloatBarrier
\subsection{Trazabilidad de decisiones con requisitos y evidencia}

Para que la justificaci\'on no quede en el plano discursivo,
se vinculan decisiones con requisitos y con evidencia emp\'irica observada
en implementaci\'on y verificaci\'on.

\begin{table}[!htbp]
\centering
\begin{tabular}{|P{1.6cm}|P{2.8cm}|P{2.9cm}|P{4.7cm}|}
\hline
\textbf{ID} & \textbf{Requisito(s) relacionado(s)} & \textbf{Evidencia principal} & \textbf{Conclusi\'on de valor} \\
\hline
D3 & RQNF-004, RQNF-008 & Separaci\'on frontend/backend en repositorio y pipelines & Mejora mantenibilidad y portabilidad \\
\hline
D8 & RQNF-009, REGN-001 & Filtros JWT y validaciones de autorizaci\'on en servicios & Refuerza control de acceso por rol/contexto \\
\hline
D11 & RQNF-001, RQNF-004 & Hist\'orico de migraciones versionadas & Reduce riesgo de deriva de esquema \\
\hline
D13 & RQNF-001, RQNF-006 & Suites unitarias/integraci\'on/E2E en verde y mutaci\'on alta & Mayor fiabilidad y menor regresi\'on \\
\hline
D16 & RQNF-008, RQNF-004 & Workflows CI/CD y despliegue automatizado en PaaS & Operaci\'on reproducible con bajo coste \\
\hline
\end{tabular}
\caption{Trazabilidad entre decisiones, requisitos y evidencia}
\label{tab:trazabilidad-decisiones}
\end{table}

\FloatBarrier
\subsection{Decisiones descartadas y por qu\'e no se adoptaron en esta fase}

\subsubsection{Autenticaci\'on federada completa en MVP}

No se adopt\'o como requisito de salida del MVP por impacto en integraci\'on,
auditor\'ia de seguridad y superficie de pruebas,
frente al valor incremental inmediato sobre el n\'ucleo funcional.

\subsubsection{2FA integral desde primera versi\'on}

Se consider\'o relevante,
pero su adopci\'on temprana pod\'ia desviar esfuerzo de procesos acad\'emicos
cr\'iticos (entregas/evaluaciones) dentro del horizonte temporal disponible.

\subsubsection{Orquestaci\'on avanzada de infraestructura}

Kubernetes y configuraciones equivalentes se descartaron
por sobrecoste operativo para un equipo reducido y un contexto de TFG,
siendo m\'as adecuado reservarlo para fases de escalado posterior.

\FloatBarrier
\subsection{Conclusiones del cap\'itulo}

La justificaci\'on detallada muestra que las decisiones del proyecto
no se basaron en preferencia tecnol\'ogica aislada,
sino en una evaluaci\'on sistem\'atica de valor, riesgo, coste y viabilidad.

El resultado es una arquitectura coherente con el ERS/DAS,
con compromisos expl\'icitos y argumentos t\'ecnicos defendibles ante tribunal,
que proyecta una hoja de ruta de evoluci\'on realista para fases futuras.
